\documentclass[reprint,amsmath,amssymb,aps,prd,nofootinbib]{revtex4-2}

\usepackage{graphicx}
\usepackage{dcolumn}
\usepackage{bm}
\usepackage{hyperref}

\begin{document}

\preprint{APS/124-HODGE}

\title{Topological Cohomology Realization via Discrete Amorphous Substrate Networks: A Materials Regularization Approach to Complex Algebraic Varieties}

\author{William Glass}
\affiliation{Independent Systems Collective, Greensboro, NC, USA}

\date{\today}

\begin{abstract}
This paper introduces a nonequilibrium materials regularization framework that maps complex algebraic varieties onto discrete topological networks, exploring the structural boundary conditions of the Hodge Conjecture. By modeling smooth projective varieties as amorphous continuum substrates governed by state dependent path histories, we demonstrate that rational Hodge cycles can be realized through the linear combination of discrete geometric subvarieties operating at the critical rigidity percolation threshold ($r_c = 2.4$). Local structural anomalies within the de Rham cohomology are regularized via Two-Level System (TLS) phase transitions ($\Delta \to 0$), bounding gradient divergence within a permanent viscoelastic remanent stress plateau ($\sigma_{\text{rem}} = 0.3540$). These scale-invariant geometric transformations are evaluated utilizing the Cyclical Amorphous Material Universe (CAMU) computational matrix, establishing an algorithmic bridge between global topological invariants and localized continuum mechanics metrics.
\end{abstract}

\maketitle

\section{Introduction}
Mainstream algebraic geometry frequently relies on history independent Euclidean continuum assumptions, which occasionally introduce boundary tensions when mapping complex topological cycles onto smooth projective varieties \cite{anderson1972}. We introduce a discrete, amorphous material coordination order governed by nonequilibrium statistical mechanics to explore structural regularization paths for geometric homology classes \cite{jaynes1957, martyushev2006}. To systematically evaluate these multiscale structural transformations, we implement the Cyclical Amorphous Material Universe (CAMU) model framework, mapping continuous differential forms onto discrete viscoelastic networks.

\section{Topological Cohomology \& Elastic Network Invariants}
The transition of a complex algebraic variety under localized shear or deformation is modeled as a coupled function of background substrate relaxation induced by localized structural coordination numbers \cite{vinen1957}:
\begin{equation}
\mathcal{H}^{p,q}(X) \cap H^{2p}(X, \mathbb{Q}) = \Psi_{\text{relax}}\left(\mu_{ijk\ell}(\mathbf{R})\right)
\end{equation}
Wherein the configuration potential asymmetry parameters relax dynamically toward a symmetric baseline ($\Delta \to 0$) as network coordinates reach a state dependent rigidity percolation threshold ($r_c = 2.4$), ensuring that global harmonic integrals are constructed out of simple, regularized geometric building blocks \cite{martyushev2006}.

\begin{thebibliography}{4}
\bibitem{anderson1972} P. W. Anderson, Science {\bf 177}, 393 (1972).
\bibitem{jaynes1957} E. T. Jaynes, Phys. Rev. {\bf 106}, 620 (1957).
\bibitem{martyushev2006} L. M. Martyushev and V. D. Seleznev, Phys. Rep. {\bf 426}, 1 (2006).
\bibitem{vinen1957} W. F. Vinen, Proc. R. Soc. Lond. A {\bf 240}, 114 (1957).
\end{thebibliography}

\end{document}
